\explainedfigure{fig_detector_retention_by_weather_intensity}{0.95}{Detector retention measured at different fog and rain intensities, and the probability that the driving model keeps a vehicle}{In 20 scenes (seeds 1000 to 1019, traffic density 0.3, stationary car, up to three vehicles placed around it) the same camera frames were passed through the teammates' vehicle detector twice: once clean and once with the weather filter at each intensity (two random weather draws per scene). Retention is the number of vehicles found with weather divided by the number found without weather. Grey dots are single scenes, the blue line is the mean over scenes with its 95\% interval, the orange dashed line is the measured retention capped at 1, which is the probability that the weather proxy keeps a vehicle in the driving model's grid. Dotted vertical lines mark the intensities used in the 200-episode weather tests.}{Fog does almost nothing up to intensity 0.6, then removes vehicles quickly: about 71\% are kept at 0.8 and 27\% at 1.0. Rain acts earlier and harder: about half of the vehicles are kept at 0.45, 16\% at 0.6 and 4\% at 1.0. Values slightly above 1 in light fog mean that the detector found a few more boxes than in the clean frame; they are treated as no loss.}{The driving model never sees the weather and never sees camera images in these tests. It receives the weather only indirectly: each vehicle is removed from its bird's-eye-view grid with the probability read off this curve. The conclusions about weather therefore hold for \emph{loss of detected vehicles at these measured rates}, from 20 scenes per point; they do not cover other perception failures.}{fig:detectorretention}
